% Generated draft: requires mathematical and source review.
\documentclass[11pt]{article}
\usepackage[margin=25mm]{geometry}
\usepackage{fontspec}
\setmainfont{Latin Modern Roman}
\usepackage{amsmath,amssymb,mathtools,unicode-math}
\setmathfont{Latin Modern Math}
\usepackage{enumitem,needspace,xurl,hyperref}
\setlength{\parindent}{0pt}
\setlength{\parskip}{5pt}
\setlength{\emergencystretch}{3em}
\newcommand{\N}{\mathbb{N}}
\newcommand{\Z}{\mathbb{Z}}
\newcommand{\Q}{\mathbb{Q}}
\newcommand{\R}{\mathbb{R}}
\newcommand{\C}{\mathbb{C}}
\newcommand{\F}{\mathbb{F}}
\newcommand{\A}{\mathbb{A}}
\newcommand{\PP}{\mathbb{P}}
\newcommand{\E}{\mathbb{E}}
\newcommand{\eps}{\varepsilon}
\newcommand{\dd}{\,\mathrm{d}}
\begin{document}
\section*{UnsolvedMath ID: 20002034}
\textbf{Source problem code:} catalogue entry
\subsection{Definitions and mathematical statement}
The source asks how \(t(g)\) and \(h(g)\) relate to other conformal invariants, without defining these symbols, signature, spin structure or function spaces. The aid interprets \(h\) as Dirac nullity and \(t\) as twistor-spinor nullity. The following fixes that interpretation with complex spinors.

On a connected closed Riemannian spin manifold \((M^n,g)\), \(n\geq3\), with fixed spin structure, let \(S\) be the complex spinor bundle, \(D_g\) its Dirac operator and
\[
 P_X\psi=\nabla_X\psi+\frac1nX\cdot D_g\psi,
 \qquad h(g)=\dim_\C\ker D_g,\quad
 t(g)=\dim_\C\ker P.
\]
Use \(X\cdot X\cdot\psi=-g(X,X)\psi\). Both nullities are conformally invariant with the standard spinor identifications and conformal weights. In this closed Riemannian category they are finite.

Let \(Y(M,[g])\) denote the Yamabe constant of the conformal class, normalized by
\[
 \inf_{u>0}
 \frac{\displaystyle\int_M
 (\tfrac{4(n-1)}{n-2}|\nabla u|^2+R_gu^2)\,\dd\mathrm{vol}_g}
 {\left(\displaystyle\int_Mu^{2n/(n-2)}\,\dd\mathrm{vol}_g\right)^{(n-2)/n}}.
\]
The aid's trichotomy is
\[
 Y>0\Longrightarrow h=0,\qquad
 Y<0\Longrightarrow t=0.
\]
If \(Y=0\), choose a scalar-flat Yamabe representative \(g_Y\). After the conformal identifications,
\[
 \ker D_{g_Y}=\ker P_{g_Y}
 =\{\psi:\nabla^{g_Y}\psi=0\},\qquad h(g)=t(g).
\]
Thus simultaneous nonzero \(h\) and \(t\) can occur only in the zero-Yamabe stratum. Zero Yamabe constant alone does not assert the existence of a nonzero parallel spinor: both nullities may be zero. The reasoning uses the closed-manifold Dirac identity \(D^2=\nabla^*\nabla+R/4\) and the twistor identity, so boundary terms or Lorentzian signs cannot be omitted from other settings.

In even dimension, decompose \(S=S^+\oplus S^-\) and let \(D^+:\Gamma(S^+)\to\Gamma(S^-)\). The aid's topological constraint is
\[
 h(g)\geq|\operatorname{ind}D^+|,
\]
since the full nullity is the sum of the two chiral kernel dimensions while the index is their difference. This is a lower bound by an index, not an upper bound on \(h\).

For contrast, on a connected globally hyperbolic Lorentzian spin manifold of dimension \(n\geq3\), the smooth spatially compact solution space of \(D\psi=0\) is infinite-dimensional. With a smooth spacelike Cauchy hypersurface \(\Sigma\), the Dirac Cauchy problem identifies it with compactly supported smooth initial spinor data on \(\Sigma\). Spatial compactness refers to support arising from compact Cauchy data and intersecting Cauchy hypersurfaces compactly. Hence even this restricted Dirac kernel has infinite complex dimension, and an unqualified smooth \(h\) is not a finite numerical invariant.

In the same Lorentzian category, twistor prolongation embeds the solution space into parallel sections of a bundle of complex rank twice the spinor rank. Connectedness therefore gives
\[
 t(g)\leq2^{\lfloor n/2\rfloor+1}.
\]
This concerns complex dimensions and full complex spinors; real or chiral conventions can change numerical counts. A Lorentzian analogue of finite Dirac nullity needs an explicitly chosen analytic domain, support condition, boundary/asymptotic condition or Fredholm formulation. The Riemannian Yamabe conclusions are not automatically transferred to this different signature.
\subsection{Short English statement}
With the aid's spinorial interpretation, Yamabe sign constrains Dirac and twistor kernels on closed Riemannian spin manifolds; at zero Yamabe sign both become parallel spinors and have equal dimension. In globally hyperbolic Lorentzian geometry, even spatially compact Dirac solutions are infinite-dimensional while twistor solutions remain finite-dimensional.
\subsection{Sources}
\begin{itemize}
\item Permanent UnsolvedMath ID \(20002034\), exact public code catalogue entry; local registry \texttt{aim-geometry-notes.json}, \texttt{notes[371]}, canonical tag problem.
\item AIM workshop ``Conformal structure in geometry, analysis, and physics,'' Problem 14: \url{https://aimath.org/WWN/confstruct/confstruct.pdf}. No PDF was accessed. The bare canonical statement does not define \(t,h\), and no proposer, year or literature entry is supplied.
\item Research attempt 1, ``Spinorial nullities, Yamabe sign, and Lorentzian finiteness,'' labelled partial\_result, partially\_solved, proved and low novelty confidence. No full references are given for the Yamabe, spinor-prolongation, index or Lorentzian Cauchy ingredients. The exact-statement verification field merely introduces the canonical record.
\end{itemize}
\end{document}
